% GENERATED FILE. Edit canonical lessons, then run npm run generate:books.
% canonical-source-hash: fnv1a64:3797d5d576f380d0
% canonical-lessons: TA-C90-laughter, TA-C90-crying, TA-C90-tears, TA-C90-alone, TA-C90-hungry

\chapter{Laughter and Tears}
\label{ch:ta-laughter-and-tears}
\hlchaptermodality{90}

\begin{chapteropening}
\textbf{By the end of this chapter:} \emph{I can talk about laughter, crying and tears, say I am alone, and say I am hungry.}

Complete the last lesson of chapter 90: I can talk about laughter, crying and tears, say I am alone, and say I am hungry.
\end{chapteropening}

\section[\texorpdfstring{\ta{சிரிப்பு} (sirippu)}{சிரிப்பு (sirippu)}]{\ta{சிரிப்பு} (sirippu) — laughter}

\label{lesson:TA-C90-laughter}



\begin{quote}
Before the new one: say the Tamil for calm, then the Tamil for shivering.
\end{quote}

\subsection*{You'll want to know: \ta{சிரிப்பு}}
\textbf{\ta{சிரிப்பு}} (\emph{sirippu}) — \textquotedblleft{}laughter\textquotedblright{}, and \textquotedblleft{}a smile\textquotedblright{}.

\textbf{\ta{சிரிப்பு} \ta{வருகிறது}} — \textquotedblleft{}I feel like laughing\textquotedblright{}, literally \textquotedblleft{}laughter is coming\textquotedblright{}.

\textbf{In the family, and next door}

\noindent
\begin{tabularx}{\linewidth}{@{}>{\raggedright\arraybackslash}X>{\raggedright\arraybackslash}X@{}}
\toprule
\textbf{} & \textbf{word} \\
\midrule
Kannada & \ka{ನಗು} (\emph{nagu}) \\
Malayalam & \ml{ചിരി} (\emph{ciri}) \\
Hindi (neighbour) & \dv{हँसी} (\emph{hãsī}) \\
English & laughter \\
\bottomrule
\end{tabularx}

\begin{grammarlens}[title={What you've built}]
Laughter, which comes to you.
\end{grammarlens}

\subsection*{Your turn}
\begin{itemize}
\raggedright
  \item \emph{Say it:} \emph{sirippu}
  \item \emph{Say it:} \emph{sirippu}, once more
  \item \emph{Say it:} say \emph{sirippu varukiṟatu}
  \item \emph{Recall:} say the Tamil for calm, then the Tamil for shivering, then say \emph{sirippu} again
\end{itemize}

\begin{tcolorbox}[breakable,colback=teal!4,colframe=teal!35!black,title={Before you move on}]
What does \ta{சிரிப்பு} mean? (\textquotedblleft{}Laughter\textquotedblright{}.) Say it once more.
\end{tcolorbox}
\section[\texorpdfstring{\ta{அழுகை} (aḻukai)}{அழுகை (aḻukai)}]{\ta{அழுகை} (aḻukai) — crying}

\label{lesson:TA-C90-crying}



\begin{quote}
Before the new one: say the Tamil for shivering, then the Tamil for laughter.
\end{quote}

\subsection*{You'll want to know: \ta{அழுகை}}
\textbf{\ta{அழுகை}} (\emph{aḻukai}) — \textquotedblleft{}crying\textquotedblright{}.

\textbf{\ta{அழுகை} \ta{வருகிறது}} — \textquotedblleft{}I feel like crying\textquotedblright{}, on the same frame as laughter.

\textbf{In the family, and next door}

\noindent
\begin{tabularx}{\linewidth}{@{}>{\raggedright\arraybackslash}X>{\raggedright\arraybackslash}X@{}}
\toprule
\textbf{} & \textbf{word} \\
\midrule
Telugu & \te{ఏడుపు} (\emph{ēḍupu}) \\
English & crying \\
\bottomrule
\end{tabularx}

\begin{grammarlens}[title={What you've built}]
Crying, in the same frame.
\end{grammarlens}

\subsection*{Your turn}
\begin{itemize}
\raggedright
  \item \emph{Say it:} \emph{aḻukai}
  \item \emph{Say it:} \emph{aḻukai}, once more
  \item \emph{Say it:} say \emph{aḻukai varukiṟatu}
  \item \emph{Recall:} say the Tamil for shivering, then the Tamil for laughter, then say \emph{aḻukai} again
\end{itemize}

\begin{tcolorbox}[breakable,colback=teal!4,colframe=teal!35!black,title={Before you move on}]
What does \ta{அழுகை} mean? (\textquotedblleft{}Crying\textquotedblright{}.) Say it once more.
\end{tcolorbox}
\section[\texorpdfstring{\ta{கண்ணீர்} (kaṇṇīr)}{கண்ணீர் (kaṇṇīr)}]{\ta{கண்ணீர்} (kaṇṇīr) — tears}

\label{lesson:TA-C90-tears}



\begin{quote}
Before the new one: say the Tamil for laughter, then the Tamil for crying.
\end{quote}

\subsection*{You'll want to know: \ta{கண்ணீர்}}
\textbf{\ta{கண்ணீர்}} (\emph{kaṇṇīr}) — \textquotedblleft{}tears\textquotedblright{}, literally \textquotedblleft{}eye-water\textquotedblright{}.

It is \textbf{\ta{கண்}}, \textquotedblleft{}eye\textquotedblright{}, and \textbf{\ta{நீர்}}, \textquotedblleft{}water\textquotedblright{}.

\textbf{In the family, and next door}

\noindent
\begin{tabularx}{\linewidth}{@{}>{\raggedright\arraybackslash}X>{\raggedright\arraybackslash}X>{\raggedright\arraybackslash}X@{}}
\toprule
\textbf{} & \textbf{word} & \textbf{same root} \\
\midrule
Kannada & \ka{ಕಣ್ಣೀರು} (\emph{kaṇṇīru}) & yes \\
Telugu & \te{కంటినీరు} (\emph{kaṇṭinīru}) & yes \\
Hindi (neighbour) & \dv{आँसू} (\emph{ā̃sū}) &  \\
English & tears &  \\
\bottomrule
\end{tabularx}

\begin{grammarlens}[title={What you've built}]
Tears, made of eye and water.
\end{grammarlens}

\subsection*{Your turn}
\begin{itemize}
\raggedright
  \item \emph{Say it:} \emph{kaṇṇīr}
  \item \emph{Say it:} \emph{kaṇṇīr}, once more
  \item \emph{Say it:} say \emph{kaṇṇīr}
  \item \emph{Recall:} say the Tamil for laughter, then the Tamil for crying, then say \emph{kaṇṇīr} again
\end{itemize}

\begin{tcolorbox}[breakable,colback=teal!4,colframe=teal!35!black,title={Before you move on}]
What does \ta{கண்ணீர்} mean? (\textquotedblleft{}Tears\textquotedblright{}.) Say it once more.
\end{tcolorbox}
\section[\texorpdfstring{\ta{தனியாக} (taṉiyāka)}{தனியாக (taṉiyāka)}]{\ta{தனியாக} (taṉiyāka) — alone}

\label{lesson:TA-C90-alone}



\begin{quote}
Before the new one: say the Tamil for crying, then the Tamil for tears.
\end{quote}

\subsection*{You'll want to know: \ta{தனியாக}}
\textbf{\ta{தனியாக}} (\emph{taṉiyāka}) — \textquotedblleft{}alone\textquotedblright{}.

\textbf{\ta{நான்} \ta{தனியாக} \ta{இருக்கிறேன்}} — \textquotedblleft{}I am alone\textquotedblright{}. \textbf{\ta{தனி}} is \textquotedblleft{}single, separate\textquotedblright{}.

\textbf{In the family, and next door}

\noindent
\begin{tabularx}{\linewidth}{@{}>{\raggedright\arraybackslash}X>{\raggedright\arraybackslash}X@{}}
\toprule
\textbf{} & \textbf{word} \\
\midrule
Telugu & \te{ఒంటరిగా} (\emph{oṇṭarigā}) \\
Hindi (neighbour) & \dv{अकेला} (\emph{akelā}) \\
English & alone \\
\bottomrule
\end{tabularx}

\begin{grammarlens}[title={What you've built}]
Alone.
\end{grammarlens}

\subsection*{Your turn}
\begin{itemize}
\raggedright
  \item \emph{Say it:} \emph{taṉiyāka}
  \item \emph{Say it:} \emph{taṉiyāka}, once more
  \item \emph{Say it:} say \emph{nāṉ taṉiyāka irukkiṟēṉ}
  \item \emph{Recall:} say the Tamil for crying, then the Tamil for tears, then say \emph{taṉiyāka} again
\end{itemize}

\begin{tcolorbox}[breakable,colback=teal!4,colframe=teal!35!black,title={Before you move on}]
What does \ta{தனியாக} mean? (\textquotedblleft{}Alone\textquotedblright{}.) Say it once more.
\end{tcolorbox}
\section[\texorpdfstring{\ta{பசிக்கிறது} (pasikkiṟatu)}{பசிக்கிறது (pasikkiṟatu)}]{\ta{பசிக்கிறது} (pasikkiṟatu) — (I) am hungry}

\label{lesson:TA-C90-hungry}



\begin{quote}
Before the new one: say the Tamil for tears, then the Tamil for alone.
\end{quote}

\subsection*{You'll want to know: \ta{பசிக்கிறது}}
\textbf{\ta{பசிக்கிறது}} (\emph{pasikkiṟatu}) — \textquotedblleft{}I am hungry\textquotedblright{}, literally \textquotedblleft{}it hungers\textquotedblright{}.

It is \textbf{\ta{பசி}}, \textquotedblleft{}hunger\textquotedblright{}, made into a verb. \textbf{\ta{எனக்குப்} \ta{பசிக்கிறது}} — \textquotedblleft{}I'm hungry\textquotedblright{}.

\begin{grammarlens}[title={What you've built}]
Laughter, crying, tears, alone, and hunger as a verb.
\end{grammarlens}

\subsection*{Your turn}
\begin{itemize}
\raggedright
  \item \emph{Say it:} \emph{pasikkiṟatu}
  \item \emph{Say it:} \emph{pasikkiṟatu}, once more
  \item \emph{Say it:} say \emph{eṉakkup pasikkiṟatu}
  \item \emph{Recall:} say the Tamil for tears, then the Tamil for alone, then say \emph{pasikkiṟatu} again
\end{itemize}

\begin{tcolorbox}[breakable,colback=teal!4,colframe=teal!35!black,title={Before you move on}]
What does \ta{பசிக்கிறது} mean? (\textquotedblleft{}(I) am hungry\textquotedblright{}.) Say it once more.
\end{tcolorbox}
